\textbf{Theorem (Cauchy–Schwarz).}  
For any real numbers \(a_{1},\dots ,a_{n}\) and \(b_{1},\dots ,b_{n}\),
\[
\Bigl(\sum_{i=1}^{n} a_i b_i\Bigr)^{2}
   \le
\Bigl(\sum_{i=1}^{n} a_i^{2}\Bigr)
\Bigl(\sum_{i=1}^{n} b_i^{2}\Bigr).
\]

\begin{proof}
\paragraph{1. Construction of a non‑negative quadratic.}  
Define
\[
f(t)=\sum_{i=1}^{n}\bigl(a_i - t b_i\bigr)^{2},\qquad t\in\mathbb{R}.
\]
Each summand is a square, hence \(f(t)\ge 0\) for every real \(t\).

\paragraph{2. Algebraic expansion.}  
Using \((x-y)^{2}=x^{2}-2xy+y^{2}\) and linearity of the sum,
\begin{align*}
f(t)
   &=\sum_{i=1}^{n}\bigl(a_i^{2}-2t a_i b_i+t^{2}b_i^{2}\bigr) \\
   &=\sum_{i=1}^{n}a_i^{2}
      -2t\sum_{i=1}^{n}a_i b_i
      +t^{2}\sum_{i=1}^{n}b_i^{2}.
\end{align*}
Thus
\begin{equation}
f(t)=\Bigl(\sum_{i=1}^{n}b_i^{2}\Bigr)t^{2}
      -2\Bigl(\sum_{i=1}^{n}a_i b_i\Bigr)t
      +\sum_{i=1}^{n}a_i^{2}. \tag{1}
\end{equation}

\paragraph{3. Applying the discriminant condition.}  
The polynomial in \((1)\) has leading coefficient
\(A=\sum_{i=1}^{n}b_i^{2}\ge 0\).  
Since \(f(t)\ge 0\) for all \(t\in\mathbb{R}\), its discriminant must be non‑positive:
\[
\Delta =\bigl(-2\sum_{i=1}^{n}a_i b_i\bigr)^{2}
        -4\bigl(\sum_{i=1}^{n}b_i^{2}\bigr)
           \bigl(\sum_{i=1}^{n}a_i^{2}\bigr)
      =4\Bigl[\bigl(\sum_{i=1}^{n}a_i b_i\bigr)^{2}
               -\bigl(\sum_{i=1}^{n}a_i^{2}\bigr)
                \bigl(\sum_{i=1}^{n}b_i^{2}\bigr)\Bigr].
\]
Thus \(\Delta\le 0\) is equivalent to
\begin{equation}
\Bigl(\sum_{i=1}^{n}a_i b_i\Bigr)^{2}
   \le
\Bigl(\sum_{i=1}^{n}a_i^{2}\Bigr)
\Bigl(\sum_{i=1}^{n}b_i^{2}\Bigr). \tag{2}
\end{equation}

\paragraph{4. Degenerate cases.}  
If \(\sum_{i=1}^{n}a_i^{2}=0\) then each \(a_i=0\) and both sides of \((2)\) equal \(0\).  
If \(\sum_{i=1}^{n}b_i^{2}=0\) then each \(b_i=0\) and again \((2)\) reduces to \(0\le 0\).  
In either situation the quadratic \(f(t)\) becomes a constant non‑negative
function, so the argument remains valid.

\paragraph{5. Equality condition.}  
Equality in \((2)\) occurs exactly when \(\Delta=0\).  For a quadratic with
non‑negative leading coefficient, \(\Delta=0\) means there is a single real
root \(\lambda\) such that
\[
f(\lambda)=\sum_{i=1}^{n}\bigl(a_i-\lambda b_i\bigr)^{2}=0.
\]
A sum of squares vanishes only when each square is zero, i.e.
\(a_i-\lambda b_i=0\) for all \(i\).  Hence
\begin{equation}
a_i=\lambda\,b_i\qquad(i=1,\dots ,n). \tag{3}
\end{equation}
Conversely, if \((3)\) holds for some \(\lambda\in\mathbb{R}\) (including
\(\lambda=0\) when one vector is the zero vector), then
\[
\sum_{i}a_i b_i=\lambda\sum_{i}b_i^{2},\qquad
\sum_{i}a_i^{2}=\lambda^{2}\sum_{i}b_i^{2},
\]
and substitution into \((2)\) yields equality.  Thus \((3)\) characterises
precisely the case of equality: the two vectors are linearly dependent.

\end{proof}